\begingroup
\renewcommand{\arraystretch}{1.45}
\fancyhead[L]{\small DETERMINATION OF ATTRACTION}
\Needspace{41\baselineskip}
\AttractionArticle{6}

The following should now be noted about the nature of this cubic equation.

I.\quad From the last term of the equation \(-aabbCC\) it is concluded that it certainly has one real root, which is either positive or, if \(C=0\), equal to zero. Let us denote this nonnegative real root by \(g\).

II.\quad Subtracting from equation 12, written as follows
\[
\begin{aligned}
&x=\frac{AAx}{aa+x}+\frac{BBx}{bb+x}+CC
\end{aligned}
\]
this equation\textperiodcentered
\[
\begin{aligned}
&g=\frac{AAg}{aa+g}+\frac{BBg}{bb+g}+CC
\end{aligned}
\]
and dividing by \(x-g\) gives a new equation containing the two remaining roots
\[
\begin{aligned}
&1=\frac{aaAA}{(aa+x)(aa+g)}+\frac{bbBB}{(bb+x)(bb+g)}
\end{aligned}
\]
which, when properly arranged and solved, gives \([14]\)
\[
\begin{aligned}
&2x=\frac{aaAA}{aa+g}+\frac{bbBB}{bb+g}-aa-bb\\[6pt]
&\qquad{}\pm\surd\!\left(\begin{gathered}\left(aa-bb-\frac{aaAA}{aa+g}+\frac{bbBB}{bb+g}\right)^2
\\[6pt]+\frac{4aabbAABB}{(aa+g)(bb+g)}\end{gathered}\right)
\end{aligned}
\]
Since the quantity under the radical sign is by its nature positive, or at least nonnegative, this expression shows that the two remaining roots are also always real.

\Needspace{28\baselineskip}
III.\quad Subtracting from one another, however, the equations written as follows
\[
\begin{aligned}
&gx=\frac{AAgx}{aa+x}+\frac{BBgx}{bb+x}+gCC\\[6pt]
&gx=\frac{AAgx}{aa+g}+\frac{BBgx}{bb+g}+xCC
\end{aligned}
\]
and dividing by \(g-x\) gives an equation containing the two remaining roots in this form:
\[
\begin{aligned}
&0=\frac{AAgx}{(aa+g)(aa+x)}+\frac{BBgx}{(bb+g)(bb+x)}+CC
\end{aligned}
\]
which clearly, if \(g\) is a positive quantity, cannot be satisfied by a positive value of \(x\). Hence we conclude that our cubic equation cannot have more than one positive root.
\SourceLeaf{351}\par
\Needspace{25\baselineskip}
IV.\quad Whenever, therefore, \(0\) is not among the roots of our equation, there will necessarily be one positive root with two negative ones. Whenever, however, \(C=0\), and hence \(0\) is one of the roots, the remaining roots will be contained in the equation
\[
\begin{aligned}
&xx-(AA+BB-aa-bb)x+aabb-aaBB-bbAA=0
\end{aligned}
\]
whence these roots will be expressed by
\[
\begin{aligned}
&\frac12(AA+BB-aa-bb)\pm\frac12\surd\!\left((AA-BB-aa+bb)^2+4AABB\right)
\end{aligned}
\]
Three cases will again have to be distinguished here.

\Needspace{20\baselineskip}
\textit{First}, if the last term \(aabb-aaBB-bbAA\) is positive (i.\ e.\ if the attracted point lies in the plane of the attracting ellipse \textit{inside} the curve), both roots, since they must be real, will have the same sign; and hence, since they cannot both be positive, they will necessarily be negative. This can also be concluded independently of what has already been proved, from the fact that the middle coefficient, which can be written as follows
\[
\begin{aligned}
&(aabb-aaBB-bbAA)\left(\frac1{aa}+\frac1{bb}\right)+\frac{bbAA}{aa}+\frac{aaBB}{bb}
\end{aligned}
\]
is clearly positive in this case.

\textit{Second}, if the last term is negative, or the attracted point is situated in the plane of the ellipse \textit{outside} the curve, one root will necessarily be positive, the other negative.

\Needspace{18\baselineskip}
\textit{Third}, however, if the last term itself were to vanish, or the attracted point were to lie on the very circumference of the ellipse, the second root too would become \(=0\), and the third
\[
\begin{aligned}
&=-\frac{bbAA}{aa}-\frac{aaBB}{bb}
\end{aligned}
\]
i.\ e.\ negative. This case, which is physically impossible and in which the attraction itself would become infinitely large, we shall exclude from our investigation, at least here.

\Needspace{46\baselineskip}
\AttractionArticle{7}

To determine the coefficients \(\gamma,\gamma',\gamma''\), from equations \(1,3,4,6,7,9,10\) we find
\SourceLeaf{352}
\[
\left.\begin{aligned}
\gamma&=\frac{1}{\surd{\bigl(1-(\frac{aA}{aa+G})^2-(\frac{bB}{bb+G})^2\bigr)}}\\[9pt]
\gamma'&=\frac{1}{\surd{\bigl((\frac{aA}{aa-G'})^2+(\frac{bB}{bb-G'})^2-1\bigr)}}\\[9pt]
\gamma''&=\frac{1}{\surd{\bigl((\frac{aA}{aa-G''})^2+(\frac{bB}{bb-G''})^2-1\bigr)}}
\end{aligned}\right\}\quad[15]
\]
From these equations, properly combined with 5, 8, 11, it also follows that:
\[
\left.\begin{aligned}
\gamma&=\surd{\frac{G}{(\frac{AG}{aa+G})^2+(\frac{BG}{bb+G})^2+CC}}\\[9pt]
\gamma'&=\surd{\frac{G'}{(\frac{AG'}{aa-G'})^2+(\frac{BG'}{bb-G'})^2+CC}}\\[9pt]
\gamma''&=\surd{\frac{G''}{(\frac{AG''}{aa-G''})^2+(\frac{BG''}{bb-G''})^2+CC}}
\end{aligned}\right\}\quad[16]
\]
These latter expressions show that none of the quantities $G,G',G''$ can be negative, if $\gamma,\gamma',\gamma''$ are to be real.

\Needspace{34\baselineskip}
In the case, therefore, where we do not have $C=0$, $G$ must necessarily be set equal to the positive root of equation $B$; and hence it is evident that $-G'$ must equal one of the negative roots, and $-G''$ the other\HistoricalNote{26}; which root we adopt for $-G'$, and which for $-G''$, will be entirely arbitrary.

Whenever $C=0$ and the attracted point is situated inside the curve, the two negative roots of equation 13 must necessarily be adopted for $-G'$ and $-G''$, and hence $G=0$ must be set. Since, however, in this case the first formula in 16 becomes indeterminate, we shall retain the first formula in 15 in its place, which gives
\[
\begin{aligned}
&\gamma=\frac{1}{\surd{\bigl(1-\frac{AA}{aa}-\frac{BB}{bb}\bigr)}}
\end{aligned}
\]
\Needspace{24\baselineskip}
Whenever, however, for $C=0$ the attracted point lies outside the ellipse. the positive root of equa\SourceLeaf{353}%
tion 13 is to be set $=G$, and either the negative root $=-G'$, and $G''=0$, or the negative root $=-G''$, and $G'=0$; the coefficient $\gamma''$ or $\gamma'$, however, we shall find by the formula
\[
\begin{aligned}
&\frac{1}{\surd{\bigl(\frac{AA}{aa}+\frac{BB}{bb}-1\bigr)}}
\end{aligned}
\]
In the case already excluded, where the attracted point were supposed to be situated on the very circumference of the ellipse, the coefficients $\gamma$ and $\gamma'$, or $\gamma$ and $\gamma''$, would become infinite, which indicates that our transformation is entirely inapplicable to this case.

\clearpage
\endgroup
